\section{Binding to the implementation}\label{sec:binding}
\status{CLOSED method and recorded findings; B1 IN-FLIGHT}
\subsection{From source text to executions}
The source3 method fixes C bytes, parses the supported C99 fragment into
a statement tree, and proves the equality between the parsed region
and that tree. Execution theorems then reason about this bound tree
under explicit state, memory and word-range contracts. This separates
the source identity from the value-level argument. It is a fragment
semantics, not a verified general C compiler.\src{C11}

For example, \code{KeygenMkgm3Program.source_bound} states
\[
 \code{C99ModularParser.region}\ 2945\ 91=\mathrm{some}\ \code{code}.
\]
The B1.02 execution work derives loop counters and pointer positions
for the first, triple and intermediate loop families of
\code{modp_NTT3_ext}. Butterfly-call observations are conclusions of
executions of the parsed bodies; assuming the intended calls would
leave the source-binding obligation unanswered. These are scoped
execution results, still subject to their memory and non-overflow
premises.\src{C11}

The generator certificate is similarly concrete. With
$p=2147355649$ and $g=1907584673$, the declarations in
\code{KeygenGeneratorOrder} prove
\begin{align*}
 g^{9216}\bmod p&=1,&g^{4608}\bmod p&\ne1,&g^{3072}\bmod p&\ne1,\\
 (g^2)^{4608}\bmod p&=1,&(g^2)^{2304}\bmod p&\ne1,&
 (g^2)^{1536}\bmod p&\ne1.
\end{align*}
The same module checks $9216=2^{10}3^2$ and $4608=2^9 3^2$.
These are the exact power tests used as the order certificates for
$9216$ and $4608$; they do not assume an NTT correctness oracle.
We quote the literal arithmetic declarations rather than silently
replacing them by a stronger source-level transform theorem.\src{C11}

\subsection{The verification found bugs}
\paragraph{F-001: an inherited C buffer-cursor defect.}
In the pinned \code{falcon_prng_get_bytes}, the available length is
computed from \code{p->ptr}, but the copy is
\code{memcpy(buf, p->buf.d, clen)} rather than a copy from
\code{p->buf.d + p->ptr}. A call beginning at a nonzero cursor can
therefore copy the buffer prefix instead of the intended stream slice.
The recorded census found the same statement in seven vendored
reference/optimized/local copies and no call sites in the examined
profiles. This is an inherited defect in the upstream lineage as
vendored here, classified LATENT; it is not evidence about every later
upstream release. The examined signing path uses the indexed
\code{get_u8}/\code{get_u64} accessors. A future use of the bulk-byte
API would activate the faulty path.\src{C12}

\paragraph{A dead REV10 parser branch.}
The initial store parser never recognized real expressions of the form
\code{gm[b + REV10[u << k]]}. Its greedy \code{pureExpr} consumed
\code{+ REV10} as part of a scalar expression, leaving a bracket behind.
The correction splits at the REV10 marker before parsing the base.
An apparently plausible row model would otherwise have described an
unbound program tree. This was a formal-model defect, not a new defect
in the C table.\src{C12}

\paragraph{Duplicated XOR and a false function call.}
The candidate \code{divSelect} included an outer XOR with \code{z},
although the compound update \code{z ^= ...} already supplied that
operation. Kernel checking rejected the claimed parse equality.
The correction uses precisely the right-hand side
\code{(z ^ z2) & -(uint32_t)((e >> i) & 1)}.
A separate parser pattern treated an opening parenthesis as a function
name and misread \code{((size_t)1 << k) - 1}; a token-name guard removed
the false call. The failed bindings and probes are retained in the
batch record.\src{C12}

\subsection{Proof engineering and current boundary}
The REV10 certificate now states
\code{rawTable = some ((List.range 1024).map bitrev10)}.
Its source-side certificate uses eleven slices, each parsing at most
eight pinned lines, and composes intermediate equalities without
re-parsing the entire table. Eight lines was the successful granularity
in the recorded experiment, not a universal limit of Lean's kernel.
The final named job is \code{keygen_rev10_cert_004}; the checkpoint
records a clean accepted run.\src{C11}

These results illustrate why a kernel proof of a model and a proof
that the model describes the source are distinct obligations. The
remaining B1 work includes source value/row laws and the eventual
emitted-material-to-fiber and law interfaces. Their completion cannot
be inferred from the closed parser or table certificates.\src{C11,C01}
\todo{B1-FINAL}{Replace this progress stub with the final source-to-law exports after delivery.}
